\hypertarget{sec:justificacion}{}%
\begin{center}
    {\Large\bfseries 2. JUSTIFICACIÓN}
\end{center}
\vspace{1em}

Este estudio busca analizar y comparar de qué manera la disponibilidad y el estado operativo de Nequi, DaviPlata y Bancolombia condicionan las actividades cotidianas e inciden en el grado de dependencia de los estudiantes de la UPTC, Sede Tunja. Dentro de la problemática ya identificada sobre la baja tenencia de efectivo y la alta vulnerabilidad en el campus, esta investigación cumple el papel de pasar de la percepción informal a la evidencia empírica. Con ello, se espera obtener conocimiento cuantitativo inédito sobre la magnitud del impacto operativo de las fallas, la clasificación de criticidad de cada plataforma y el nivel real de preparación transaccional de los estudiantes, información que hoy no existe a nivel local.

Es necesario investigarlo ahora porque la adopción de billeteras digitales en Colombia se ha masificado aceleradamente y las caídas o fallas masivas de estas aplicaciones son un fenómeno real y recurrente en el país. En el campus de la UPTC Sede Tunja, los estudiantes realizan múltiples pagos de bajo monto cada semana (alimentación, fotocopias y transporte) bajo horarios fijos.

El estudio sirve para que los estudiantes reconozcan su grado de dependencia y la importancia de contar con alternativas de pago en su rutina académica. A su vez, para la Ingeniería de Sistemas representa un aporte relevante al evaluar las consecuencias de la indisponibilidad de software en la vida real, subrayando la necesidad de desarrollar sistemas de información transaccionales con alta disponibilidad, tolerancia a fallas y arquitecturas resilientes ante la alta demanda.

La investigación se aborda mediante un enfoque descriptivo y comparativo sustentado en un instrumento estandarizado de recolección de datos, el cual resulta más conveniente que una simple encuesta de percepción aislada porque permite relacionar directamente eventos de falla operacionales con consecuencias medibles en la cotidianidad de los usuarios. Además, este diseño metodológico deja una matriz de evaluación replicable que servirá de base para futuras investigaciones sobre usabilidad, disponibilidad y dependencia de servicios digitales en otras instituciones de educación superior de la región.
